\documentclass[a4paper]{article} 
\addtolength{\hoffset}{-2.25cm}
\addtolength{\textwidth}{4.5cm}
\addtolength{\voffset}{-3.25cm}
\addtolength{\textheight}{5cm}
\setlength{\parskip}{0pt}
\setlength{\parindent}{0in}

%----------------------------------------------------------------------------------------
%	PACKAGES AND OTHER DOCUMENT CONFIGURATIONS
%----------------------------------------------------------------------------------------

\usepackage{blindtext} % Package to generate dummy text
\usepackage{charter} % Use the Charter font
\usepackage[utf8]{inputenc} % Use UTF-8 encoding
\usepackage{microtype} % Slightly tweak font spacing for aesthetics
\usepackage[english, spanish, es-nodecimaldot]{babel} % Language hyphenation and typographical rules
\usepackage{amsthm, amsmath, amssymb} % Mathematical typesetting
\usepackage{float} % Improved interface for floating objects
\usepackage[final, colorlinks = true, 
            linkcolor = black, 
            citecolor = black]{hyperref} % For hyperlinks in the PDF
\usepackage{graphicx, multicol} % Enhanced support for graphics
\usepackage{subcaption} % For subfigures
\usepackage{xcolor} % Driver-independent color extensions
\usepackage{marvosym, wasysym} % More symbols
\usepackage{rotating} % Rotation tools
\usepackage{censor} % Facilities for controlling restricted text
\usepackage{pseudocode} % Environment for specifying algorithms in a natural way
\usepackage{booktabs} % Enhances quality of tables
\usepackage{tikz-qtree} % Easy tree drawing tool
\tikzset{every tree node/.style={align=center,anchor=north},
         level distance=2cm} % Configuration for q-trees
\usepackage{csquotes} % Context sensitive quotation facilities
\usepackage[yyyymmdd]{datetime} % Uses YEAR-MONTH-DAY format for dates
\renewcommand{\dateseparator}{-} % Sets dateseparator to '-'
\usepackage{fancyhdr} % Headers and footers
\pagestyle{fancy} % All pages have headers and footers
\fancyhead{}\renewcommand{\headrulewidth}{0pt} % Blank out the default header
\fancyfoot[L]{} % Custom footer text
\fancyfoot[C]{} % Custom footer text
\fancyfoot[R]{\thepage} % Custom footer text
\newcommand{\note}[1]{\marginpar{\scriptsize \textcolor{red}{#1}}} % Enables comments in red on margin
\DeclareMathOperator*{\plim}{plim}
\usepackage[most]{tcolorbox}
\usepackage{cancel}
\usepackage{adjustbox}
\addto\captionsspanish{
\def\listtablename{\'Indice de tablas}%
\def\tablename{Tabla}}
\usepackage{xcolor}
\usepackage{multirow}
\usepackage{listings}
\usepackage{bbm}
\usepackage{bm}
\usepackage{dsfont}
\usepackage{listings}
\usepackage{xcolor}
\usepackage{hyperref}
%Añadí estos dos: 
\usepackage{changepage}
\usepackage{csquotes}
\usepackage[colorlinks]{hyperref}




\definecolor{myblue}{RGB}{0,163,243}
\definecolor{moradito}{RGB}{63,1,143}
\begin{document}

%-------------------------------
%	TITULO
%-------------------------------

\fancyhead[C]{}
\hrule \medskip 
\begin{minipage}{0.295\textwidth} 
\raggedright
\textbf{Profesor:} María del Pilar López Uribe\\
\vspace{2mm}
\textbf{Integrantes:} \\
Edmundo Arias - 202110688\\
Kelly Cadena - \\
Juan José Gutiérrez - \\
Nicolás Martínez - 202122948\\



\end{minipage}
\begin{minipage}{0.4\textwidth} 
\centering 
\huge 
Entrega 3\\ 
\vspace{2mm}
\normalsize 
HEC – 2024 1\\ 
Abril 01, 2024
\end{minipage}
\begin{minipage}{0.295\textwidth} 
\begin{figure}[H]
\raggedleft
\includegraphics[scale=0.3]{HEC 2024-1/Format/Style/uniandes.pdf}
\end{figure}
\hfill
\end{minipage}
\medskip\hrule 
\bigskip

\section{Título \& Tema}

\subsection{Justificación del cambio del tema}

\section{Resumen}
\subsection{Abstract}
El siguiente artículo estudia los efectos de la violencia en las posturas políticas de los municipios colombianos, enfocándose en el impacto que tuvo la violencia entre las décadas de 1970 y 1990, un periodo crítico en la historia de Colombia. Investigamos cómo esta época de acentuada violencia modificó las decisiones políticas a nivel municipal, específicamente en términos del comportamiento electoral, como se refleja en el número de votos recibidos por los partidos políticos en las elecciones locales al Concejo. Hipotetizamos que la violencia benefició los patrones de voto hacia ciertos partidos políticos, dadas sus distintas posturas ante los grupos armados ilegales. Analizamos cómo los eventos violentos podrían haber hecho que los votantes se mostraran más reacios o predispuestos a apoyar a un partido político específico. Nuestros hallazgos sugieren una relación causal: mediante una estrategia de identificación de estudio de eventos, encontramos evidencia que sugiere que los municipios que experimentaron ataques violentos durante el periodo estudiado tendieron a votar en mayor proporción por el Partido Conservador, sin que ello afectara significativamente el apoyo hacia otros partidos o el voto en blanco. Además, investigamos la persistencia de estos efectos, descubriendo una disminución significativa del apoyo al Partido Liberal en años posteriores (1978-2007). Por último, analizamos efectos de persistencia sobre otras variables de interés como confianza hacia el gobierno, gasto público local, inversión total en el municipio, y demás, utilizando el Panel Municipal CEDE. Nuestros hallazgos se suman a la literatura que estudia los determinantes de las posturas políticas, destacando la violencia como un factor crucial en la configuración de las decisiones políticas de los actores involucrados.

\subsection{Contribución}

\section{Datos}

\section{Metodologías}

\subsection{Introducción \& Especificación}
\begin{adjustwidth}{0.5}{0.5}


Nuestro estudio tendrá un énfasis causal, es decir, su objetivo final será estimar el efecto causal que tuvieron los distintos ataques violentos transcurridos entre 1972 y 1990 sobre la votación de los municipios en las elecciones al concejo. En particular, centraremos nuestra atención hacia los votos que recibió el Partido Conservador para cada año electoral que se encuentre en el periodo de estudio. Para cuantificar el efecto, utilizaremos a la violencia como una variable dicótoma que defina el tratamiento para cada municipio. En ese orden de ideas, los municipios afectados por la violencia durante los años de estudio recibirán el valor de uno y cero de lo contrario. Claramente, la presencia de violencia en los municipios no es un tratamiento homogéneo en el tiempo: incluso los municipios de un mismo departamento pudieron verse afectados por hechos de violencia guerrillera o paramilitar en distintos periodos de tiempo, por lo que un modelo de Diferencias en Diferencias (DiD) tradicional no sería la mejor herramienta para identificar el efecto causal. Es por esta razón que utilizaremos el  Estudio de Eventos como estrategia de identificación; este modelo es una generalización de (DiD) en el que las unidades de estudio pueden recibir el tratamiento en periodos distintos. Con el Estudio de Eventos, aprovecharemos la riqueza de los datos panel para encontrar el efecto que de la violencia en la cantidad de votos por el concejo que recibió el Partido Conservador. \\


Planteamos el siguiente modelo a estimar por Mínimos Cuadrados Ordinarios:
\begin{equation}
ln(v_{it})=\phi_t+\nu_i + \sum_{j=-18}^{-2}\tau_{-j}D_i\mathds{1}[j=t-t_0] + \sum_{j=2}^{18}\tau_{j}D_i\mathds{1}[j=t-t_0] + \epsilon_{it} 
\end{equation}

\begin{itemize}
\item $ln(v_{it})$ representa el logaritmo de la cantidad de votos por el Partido Conservador del municipio $i$ en el periodo $t$
\item \( \phi_t \) , \(  \nu_i\) capturan efectos fijos de período $t$ y municipio $i$.
\item \( D_{i} \) es una dicótoma que toma el valor de uno si el municipio $i$ recibió violencia algún período
\item \(\tau_{-j}\) son los \textit{leads} del efecto: el efecto causal de la violencia antes de que un municipio fuera atacado. En últimas, estos estimadores nos permiten validar el supuesto de identificación de \textit{tendencias paralelas}.
\item \(\tau_{j}\) son los \textit{lags} del efecto que caracteriza los efectos dinámicos del tratamiento: nos permiten ver la magnitud y significancia estadística del efecto de la violencia sobre los patrones de votos de un municipio a lo largo de los años.



\end{adjustwidth}

\subsection{Ventajas \& Desventajas}
Existen varias ventajas al emplear la metodología de estudio de eventos. La mayor ventaja es que se presta para presentar evidencia sugestiva del supuesto de identificación: tendencias paralelas. Es decir, podemos estimar los parámetros que indican si hay diferencias sistemáticas en la forma en cómo los patrones de votación de los municipios evolucionan en el tiempo para de tal manera identificar si estos cambian o no de manera paralela. Asimismo, a través de los mismos estimadores es posible evidenciar si el tratamiento (en nuestro caso los ataques violentos) sí tienen o no un efecto sobre los patrones de votos de los municipios afectados puesto que podemos proyectar las estimaciones en el tiempo para ver sus efectos dinámicos. Una gráfica de este estilo es presentada en la sección de \textit{Resultados}. \\

Otra ventaja no menospreciable es que esta metodología es ampliamente utilizada en la literatura económica y en la frontera de la investigación. De hecho, la estrategia de diferencias en diferencias, y su extensión de tratamiento escalonado que nosotros empleamos (conocida popularmente como estudio de eventos) es la metodología empírica más utilizada en investigadores que buscan estimar un efecto causal \cite{Angrist2022EmpiricalEffect}. Por ende, nuestros resultados se basan en métodos transparantes, robustos y ampliamente conocidos por los científicos sociales. \\

No obstante, nuestra metodología no está eximida de sus limitaciones y desvantajas, la principal siendo que los estudios de eventos son demandantes en términos de datos. Específicamente, se requieren de datos panel con varias unidades y años tanto antes como después del tratamiento para estimar el efecto causal y para proveer evidencia sugestiva del supuesto de identificación. En nuestro caso, al tratarse de datos históricos de Colombia, tuvimos que realizar un arduo trabajo para conseguir, recopilar y en últimas armar una base de datos tipo panel en la cual emplear esta metodología. Ahora bien, al tratarse de datos históricos y del hecho que hubo un grado considerable de trabajo manual en la construcción de los datos, pueden también haber errores de medición en las variables o que hayan incompletitudes para algunos municipios en ciertos años. \\

Una limitación secundaria propia de la metodología es la gran `explosión' de literatura que ha habido en este tema en particular en los últimos cinco años. En últimas, al ser una metodología ubicada en la frontera del conocimiento se presta para que hayan errores los cuales la literatura han no logra resolver. Un par de ejemplos siendo el problema de los \textit{negative weights} que se les da a los diferentes grupos de unidades con respecto a cuándo y qué tan expuestos están al tratamiento, y el problema que aún no se ha decidido definitivamente cuál es el mejor estimador para utilizar.


\section{Resultados}
\subsection{Estadísticas Descriptivas}


\subsection{Resultado Principal}
Tras estimar la ecuación (1) por MCO y guardando los coeficientes estimados para luego graficarlos llegamos al siguiente resultado:
\begin{figure}[H]
    \centering
    \caption{Efecto Causal de la Violencia sobre los Patrones de Voto de Municipios (1972 - 1990)}
    \includegraphics[width=0.5\linewidth]{HEC 2024-1//Entregas//Results/TWFE_1.pdf}
    \label{fig:Main}
\end{figure}

De la \textit{Figura 1} podemos evidenciar un claro cambio en los patrones de votos de los municipios que sufrieron un ataque violento (a favor del Partido Conservador). Es evidente que previo al tratamiento (un ataque violento) los municipios que sufrieron ataques y aquellos que no, no tenían diferencias sistemáticas en sus patrones de voto, razón por la cual todos los estimadores antes del año cero son estadísticamente equivalentes a cero. No obstante, observamos que una vez un municipio sufre un ataque sus patrones de voto a favor del Partido Conservador aumentan en los años subsecuentes. Específicamente, si un municipio sufrió de un ataque violento entre 1972 y 1990 sus votos al concejo a favor del Partido Conservador aumentarán entre un 25\% y 50\% \textit{ceteris paribus}. \\

De este anterior resultado queremos hacer una aclaración. Si bien una objección posible sería que se resalte el hecho que las estimaciones son ruidodas en ciertos años, se debe tener en cuenta que contamos con datos históricos, recopilados y construídos propiamente por los autores. Aún así, enfatizamos la validez de este resultado ya que como demostramos en la \textit{Figura 2} del Apéndice 7.2 no hay cambios estadísticamente significativos en los patrones de votos para los demás partidos (ni del voto en blanco), es decir, el Partido Conservador es el único partido que debido a la violencia observa un aumento sustantivo de sus votos en el período estudiado. \\

Ahora bien, en el contexto más amplio del artículo este resultado le añade un grado de validez empírica considerable a nuestro marco teórico. Hemos demostrado empiricamente que existe una relacion causal entre los ataques violentos sobre una población y las actitudes politicas que dicha poblacion refleja mediante sus patrones de votos en años subsecuentes. \textbf{\textcolor{blue}{Terminar una vez tengamos secciones 1 y 2}}

\subsection{Persistencia}
Una de las extensiones que queremos estudiar a la luz de los resultados anteriores, es cómo un aumento a favor de los votos hacia el Partido Conservador a causa de la violencia tiene repercusiones sobre otras variables de interés en períods futuros. Nos enfocaremos en dos posibles canales de persistencia: patrones de votos futuros, y diferencias significativas entre municipios afectados y no afectados por la violencia (utilizando variables de resultado extraídas del Panel CEDE).

\subsubsection{Patrones de Votos Futuros}
Para analizar si el cambio en los patrones de votos de los municipios a favor del Partido Conservador a causa de ataques violentos tuvo efectos persistentes en el tiempo ampliamos el período de estudio para 1978 y 2007. Esta ampliación la hacemos por varias razones, la primera siendo que obtenemos una mayor variación de los datos de violencia por municipio utilizando datos más recientes. Los datos de violencia provienen de reportes realizados por personas y es probable que en años anteriores (1972-1990) personas que si bien fueron afectadas simplemente no reportaran sus casos. No obstante, evidenciamos que hay un número de reportes de violencia, y por ende, mayores casos de violencia a nivel municipal entre 1990 y 2007. La segunda razón por la que ampliamos el panel es porque ver años subsecuentes a los estudiados (1991 - 2007) nos permitirá identificar si hay o no persistencia y en el caso afirmativo, la significancia económica de dichos efectos duraderos. \\

Con base en lo anterior, estimamos exactamente la misma ecuación (1) por MCO y llegamos al siguiente resultado:

\begin{figure}[H]
    \centering
    \caption{Persistencia de la Violencia sobre el Voto Liberal (1990 - 2007)}
    \includegraphics[width=0.5\linewidth]{HEC 2024-1//Entregas//Figures/TWFE_Ext1.pdf}
    \label{fig:persistence1}
\end{figure}

Una vez más evidenciamos que si bien las estimaciones son ruidosas, hay un claro bajón en la cantidad de votos que recibe el Partido Liberal en elecciones de Concejo entre 1990 y 2007 debido a la violencia. Específicamente, el hecho de que un municipio haya experimentado un ataque violento durante este período causa una disminución de entre 25\% y 50\% de los votos obtenidos por el Partido Liberal \textit{ceteris paribus}. Al igual que en los resultados de la \textit{Sección 5.2}, el Partido Liberal es el único partido que exhibe una caida estadísticamente significante en los años subsecuentes. \\

Este resultado anterior presenta indicios de persistencia de los efectos de la violencia, a favor del voto conservador en el corto y mediano plazo, y en contra del voto liberal en el largo plazo. En últimas, para entender estos resultados debemos tener presenta el clima político por el cual pasaba Colombia durante el período estudiado. Este fue uno de violencia, conflicto interno y marginalización política que se manifestó en movilizaciones bélicas y en el surgimiento de grupos al margen de la ley. Esta polarización política que azotaba al país, y que ha estado presente desde su concepción juega un factor fundamental en la configuración de las posturas y actitudes políticas de las personas más expuestas a la violencia.

\subsubsection{Otras Variables}
En esta sección del artículo buscaremos evaluar si la violencia que experimentó un municipio durante la década de los 1970 hasta 1990 tiene repercusiones de largo plazo hoy en día. Para ello, nos apoyaremos del Panel Municipal CEDE, el cual contiene información valiosa de distintos indicadores y variables de interés como por ejemplo: confianza hacia el Estado, nivel de gasto público, inversión total hacia el municipio, etc. Así, buscaremos establecer vínculos históricos entre la violencia que un municipio experimentó hace cincuenta años y distintas variables de interés en el presente. Para esto tenemos dos opciones, podemos reutilizar los datos originales de violencia de los cuales obtenemos el resultado principal, y realizar proyecciones locales \cite{Jorda2005EstimationProjections} para de tal manera analizar cómo la violencia histórica influye y tienen efectos que persisten hasta hoy en día. La segunda opción sería únicamente utilizar las variables indicadoras de violencia (y restringir a aquellas que indican violencia sobre nuestros años de estudio) del Panel CEDE para comparar las variables de interés sobre aquellos municipios afectados y aquellos no. No obstante, esta segunda opción sería exclusivamente descriptiva y no podríamos atribuirle una interpretación causal a los resultados.


\section{Referencias}
\bibliographystyle{apalike}
\bibliography{references}



\section{Apéndice}
\subsection{Datos}


\subsection{Resultados}
\begin{figure}[H]
    \centering
    \caption{Efecto de la Violencia sobre Otros Patrones de Votos}
    % Top left figure
    \begin{subfigure}[b]{0.45\linewidth}
        \includegraphics[width=\linewidth]{HEC 2024-1/Entregas/Results/TWFE_2.pdf}
        \caption{Voto Liberal}
        \label{fig:sub1}
    \end{subfigure}
    \hfill % This adds space between your two top figures
    % Top right figure
    \begin{subfigure}[b]{0.45\linewidth}
        \includegraphics[width=\linewidth]{HEC 2024-1/Entregas/Results/TWFE_3.pdf}
        \caption{Voto en Blanco}
        \label{fig:sub2}
    \end{subfigure}
    
    % Bottom center figure
    \begin{subfigure}[b]{0.6\linewidth} % Adjust the linewidth as needed for sizing
        \centering
        \includegraphics[width=\linewidth]{HEC 2024-1/Entregas/Results/TWFE_4.pdf}
        \caption{Voto Total}
        \label{fig:sub3}
    \end{subfigure}
    \label{fig:robustness1}
\end{figure}


\end{document}

